% Source: https://artofproblemsolving.com/community/c3h1273804_box_help
% Stable identity: aops-topic-1273804-q01
Wrapping paper covers surface area, which scales with the square of the linear dimensions. Halving every dimension multiplies the surface area by $\left(\dfrac{1}{2}\right)^2=\dfrac{1}{4}$.

\begin{asy}
usepackage("amsmath");
size(5cm);
defaultpen(fontsize(10pt));
draw(box((0,0),(4,2)),blue); draw((0,2)--(1,2.6)--(5,2.6)--(4,2),blue); draw((4,0)--(5,0.6)--(5,2.6),blue);
draw(box((6.5,0),(8.5,1)),red); draw((6.5,1)--(7,1.3)--(9,1.3)--(8.5,1),red); draw((8.5,0)--(9,0.3)--(9,1.3),red);
label("2 rolls",(2.5,0),S,blue); label("$\dfrac{1}{2}$ roll",(7.75,0),S,red);
label("\mbox{}",(min(currentpicture,user=true)+max(currentpicture,user=true))/2,black); // keeps the text after the figure black when LaTeX typesets the labels inline
\end{asy}

So the smaller box needs $2\cdot\dfrac{1}{4}=\boxed{\dfrac{1}{2}}$ roll.
